\documentclass[11pt]{article}
\usepackage[margin=2.2cm]{geometry}
\usepackage{fontspec}
\setmainfont{STIX Two Text}
\usepackage{unicode-math}
\setmathfont{STIX Two Math}
\usepackage{graphicx,booktabs,enumitem,hyperref,float,amsmath}
\setlist{nosep,leftmargin=1.4em}
\setlength{\parskip}{5pt}\setlength{\parindent}{0pt}
\newcommand{\step}[1]{\textbf{Step #1.}\ }
\newcommand{\chk}[1]{\par\smallskip\noindent\textbf{Check.}\ #1\par}
\hypersetup{hidelinks}

\setlength{\emergencystretch}{2em}

\title{Derivations, part 2 of 3:\\ the vector (Richards--Wolf) focal field of one lens}
\date{18 November 2026}
\author{}
\begin{document}
\maketitle
\vspace{-2.5em}

Part 2 derives the field of an aplanatic lens with its polarization, the model of Romallosa, Bantang \& Saloma
(2003) [Rom03], which the current work uses at high NA. Part 1 derived the scalar field; Part 3 combines two beams.
The numbers of each check are printed by \texttt{scripts/derivation\_checks.py}.

\section*{Symbols}
\begin{center}\small\begin{tabular}{ll}\toprule
$(r, \psi, z)$ & cylindrical coordinates of the point $P$; $\psi$ is measured from the input polarization\\
$(\vartheta, \varphi)$ & polar and azimuthal angles of the direction of a ray converging on the focus\\
$\hat{\mathbf k}$ & direction of that ray, $(\sin\vartheta\cos\varphi,\ \sin\vartheta\sin\varphi,\ \cos\vartheta)$\\
$\hat{\mathbf e}$ & unit field of that ray after the lens\\
$h$, $f$ & height of the ray in the pupil; focal length\\
$A(\vartheta)$ & profile of the input beam (1 for uniform input)\\
$I_0, I_1, I_2$ & the three integrals of the field, Section 3\\
$t$ & $\cos\vartheta$\\\bottomrule
\end{tabular}\end{center}

\section{The aplanatic factor}
\step{1} An aplanatic lens obeys the sine condition: the ray that enters the pupil at height $h$ leaves toward the
focus at the angle $\vartheta$ with
\begin{equation}h = f\sin\vartheta .\end{equation}

\step{2} The power that enters through the ring of the pupil between $h$ and $h + dh$ is
$|E_{\rm in}|^2\,2\pi h\,dh$.

\step{3} The same power crosses the band of the reference sphere (radius $f$, centred on the focus) between
$\vartheta$ and $\vartheta + d\vartheta$, whose area is $2\pi f^2\sin\vartheta\,d\vartheta$. With $E_{\rm s}$ the
field on the sphere,
\begin{equation}|E_{\rm in}|^2\,2\pi h\,dh = |E_{\rm s}|^2\,2\pi f^2\sin\vartheta\,d\vartheta .\end{equation}

\step{4} From Step 1, $dh = f\cos\vartheta\,d\vartheta$. Substituting $h$ and $dh$ on the left,
\begin{equation}|E_{\rm in}|^2\,2\pi f\sin\vartheta\,f\cos\vartheta\,d\vartheta
= |E_{\rm s}|^2\,2\pi f^2\sin\vartheta\,d\vartheta .\end{equation}

\step{5} The factor $2\pi f^2\sin\vartheta\,d\vartheta$ appears on both sides and cancels:
\begin{equation}|E_{\rm s}|^2 = |E_{\rm in}|^2\cos\vartheta,\qquad |E_{\rm s}| = A(\vartheta)\sqrt{\cos\vartheta} .\end{equation}
The rays at large $\vartheta$ come from a wider ring of the pupil but spread over a still wider band of the sphere,
so their amplitude falls as $\sqrt{\cos\vartheta}$.

\section{The field of one ray after the lens}
\step{1} The input is polarized along $x$: $\mathbf E_{\rm in} = \hat{\mathbf x}$. At the azimuth $\varphi$ I split
it into the radial and the azimuthal unit vectors of the pupil,
$\hat{\mathbf e}_\rho = (\cos\varphi, \sin\varphi, 0)$ and $\hat{\mathbf e}_\varphi = (-\sin\varphi, \cos\varphi, 0)$:
\begin{equation}\hat{\mathbf x} = \cos\varphi\,\hat{\mathbf e}_\rho - \sin\varphi\,\hat{\mathbf e}_\varphi .\end{equation}
(Check: $\cos\varphi(\cos\varphi, \sin\varphi) - \sin\varphi(-\sin\varphi, \cos\varphi) = (1, 0)$.)

\step{2} The lens bends the ray within its meridional plane, the plane that holds the axis and the ray. The azimuthal
component is perpendicular to that plane and does not change. The radial component turns with the ray and becomes the
meridional unit vector
\begin{equation}\hat{\mathbf e}_\vartheta = (\cos\vartheta\cos\varphi,\ \cos\vartheta\sin\varphi,\ -\sin\vartheta) ,\end{equation}
which is perpendicular to $\hat{\mathbf k}$ and equals $\hat{\mathbf e}_\rho$ when $\vartheta = 0$. The ray that arrives
with the direction $(\vartheta, \varphi)$ leaves the pupil on the opposite side, at the azimuth $\varphi + \pi$; there
both unit vectors of the pupil and both coefficients of Step 1 change sign, so the result below is the same.

\step{3} The field of the ray on the reference sphere is then
\begin{equation}\hat{\mathbf e} = \cos\varphi\,\hat{\mathbf e}_\vartheta - \sin\varphi\,\hat{\mathbf e}_\varphi
= \big(\cos\vartheta\cos^2\!\varphi + \sin^2\!\varphi,\ \ (\cos\vartheta - 1)\sin\varphi\cos\varphi,\ \
-\sin\vartheta\cos\varphi\big) .\end{equation}

\step{4} With $\cos^2\!\varphi = \frac12(1 + \cos2\varphi)$, $\sin^2\!\varphi = \frac12(1 - \cos2\varphi)$ and
$\sin\varphi\cos\varphi = \frac12\sin2\varphi$:
\begin{align}
e_x &= \tfrac12(1 + \cos\vartheta) - \tfrac12(1 - \cos\vartheta)\cos2\varphi ,\\
e_y &= -\tfrac12(1 - \cos\vartheta)\sin2\varphi ,\\
e_z &= -\sin\vartheta\cos\varphi .
\end{align}
Each component is now a sum of terms $\cos m\varphi$ or $\sin m\varphi$ with $m = 0$, 1 or 2.

\chk{$\hat{\mathbf k}\cdot\hat{\mathbf e} = \sin\vartheta\cos\varphi\,[\cos\vartheta(\cos^2\!\varphi + \sin^2\!\varphi) +
\sin^2\!\varphi - \sin^2\!\varphi - \cos\vartheta] = 0$: the field of each ray is transverse, as a plane wave must
be. On a grid of 141 by 181 directions the largest $|\hat{\mathbf k}\cdot\hat{\mathbf e}|$ is $1.7\times10^{-16}$ and
$|\hat{\mathbf e}|$ differs from 1 by at most $3.3\times10^{-16}$.}

\section{The field near the focus}
\step{1} As in Part 1, Section 1, the field is the sum of the plane waves, now each with its vector
$\hat{\mathbf e}$ and the amplitude of Section 1:
\begin{equation}\mathbf E(r,\psi,z) = C\int_0^\alpha\!\!\int_0^{2\pi}\hat{\mathbf e}(\vartheta,\varphi)\,
A(\vartheta)\sqrt{\cos\vartheta}\;e^{ik(r\sin\vartheta\cos(\varphi-\psi) + z\cos\vartheta)}\,
\sin\vartheta\,d\varphi\,d\vartheta ,\end{equation}
with one constant $C$ for the three components.

\step{2, the integrals over $\varphi$} I write $\varphi = \psi + \chi$. Then
$\cos m\varphi = \cos m\psi\cos m\chi - \sin m\psi\sin m\chi$. The term with $\sin m\chi$ integrates to zero over a
period, because $\sin m\chi$ is odd in $\chi$ and $e^{iv\cos\chi}$ is even. The other term is Bessel's integral,
$\int_0^{2\pi}\cos m\chi\,e^{iv\cos\chi}\,d\chi = 2\pi\,i^mJ_m(v)$. In the same way for $\sin m\varphi$:
\begin{equation}\int_0^{2\pi}\!\cos m\varphi\;e^{iv\cos(\varphi-\psi)}d\varphi = 2\pi\,i^mJ_m(v)\cos m\psi,\qquad
\int_0^{2\pi}\!\sin m\varphi\;e^{iv\cos(\varphi-\psi)}d\varphi = 2\pi\,i^mJ_m(v)\sin m\psi ,\end{equation}
with $v = kr\sin\vartheta$, $i^0 = 1$, $i^1 = i$ and $i^2 = -1$.

\step{3, the integrals over $\vartheta$} What remains of each term is one of three integrals over $\vartheta$:
\begin{align}
I_0 &= \int_0^\alpha A\sqrt{\cos\vartheta}\,\sin\vartheta\,(1 + \cos\vartheta)\,J_0(kr\sin\vartheta)\,e^{ikz\cos\vartheta}\,d\vartheta ,\\
I_1 &= \int_0^\alpha A\sqrt{\cos\vartheta}\,\sin^2\!\vartheta\,J_1(kr\sin\vartheta)\,e^{ikz\cos\vartheta}\,d\vartheta ,\\
I_2 &= \int_0^\alpha A\sqrt{\cos\vartheta}\,\sin\vartheta\,(1 - \cos\vartheta)\,J_2(kr\sin\vartheta)\,e^{ikz\cos\vartheta}\,d\vartheta .
\end{align}
These are Eqs.~5 to 7 of [Rom03], written there with $u = k\sin^2\!\alpha\,z$ and $v = k\sin\alpha\,r$.

\step{4, $E_x$} The term $\frac12(1 + \cos\vartheta)$ has $m = 0$ and gives $C\cdot\frac12\cdot2\pi I_0 = C\pi I_0$.
The term $-\frac12(1 - \cos\vartheta)\cos2\varphi$ has $m = 2$; its minus sign and the $i^2 = -1$ cancel, and it
gives $C\pi I_2\cos2\psi$:
\begin{equation}E_x = C\pi\,(I_0 + I_2\cos2\psi) .\end{equation}

\step{5, $E_y$} The term $-\frac12(1 - \cos\vartheta)\sin2\varphi$ has $m = 2$ and gives, in the same way,
\begin{equation}E_y = C\pi\,I_2\sin2\psi .\end{equation}

\step{6, $E_z$} The term $-\sin\vartheta\cos\varphi$ has $m = 1$. Its $\sin\vartheta$ and the $\sin\vartheta$ of the
solid angle make the $\sin^2\!\vartheta$ of $I_1$, and $i^1 = i$:
\begin{equation}E_z = C\,(-1)\,(2\pi i)\,I_1\cos\psi = -2i\,C\pi\,I_1\cos\psi .\end{equation}

\step{7, the constant} Richards \& Wolf, and [Rom03], fix the constant so that $E_x = -i(I_0 + I_2\cos2\psi)$, that is
$C\pi = -i$. The same constant multiplies $E_y$ and $E_z$:
\begin{equation}E_x = -i\,(I_0 + I_2\cos2\psi),\qquad E_y = -i\,I_2\sin2\psi,\qquad
E_z = -2i\,(-i)\,I_1\cos\psi = -2\,I_1\cos\psi ,\label{eq:rw}\end{equation}
since $-2i\cdot(-i) = 2i^2 = -2$. The factor of $E_z$ is real: $E_z$ is a quarter period out of phase with $E_x$. The
two-beam code had $-2i$ until 2026-11-10, which used the $i$ of Step 6 and left out the $-i$ of the constant.

\chk{Eqs.~2 to 4 of [Rom03] as printed read $E_x = -j(I_0 + 2I_2\cos\phi)$ and $E_y = -jI_2\sin\phi$, where
Eq.~(\ref{eq:rw}) has $\cos2\psi$ and $\sin2\psi$; $E_z = -2I_1\cos\phi$ is as above. The width that [Rom03] reports
decides which form their computations used. At NA 0.8, $n = 1.3$ and 750 nm, the FWHM across the polarization is
472.8 nm with Eq.~(\ref{eq:rw}), $|I_0 - I_2|^2$, which is the value of [Rom03]; the printed form gives 482.1 nm there.
Along the polarization Eq.~(\ref{eq:rw}) gives 525.1 nm and the printed form 536.8 nm. The printed $\cos\phi$ and
$\sin\phi$ are therefore a slip of typesetting, and the engine uses Eq.~(\ref{eq:rw}). A term $\cos\phi$ would also make
$E_x$ change under $r \to -r$ in the focal plane, which the symmetric lens does not allow.}

\step{8, the same integral in all planes} From Eq.~(\ref{eq:rw}), the mean of $|\mathbf E|^2$ over the azimuth
$\psi$ is
\begin{equation}\langle|\mathbf E|^2\rangle_\psi = |I_0|^2 + |I_2|^2 + 2|I_1|^2 ,\end{equation}
since the mean of $\cos2\psi$ is 0 and the means of $\cos^2 2\psi$, $\sin^2 2\psi$ and $\cos^2\psi$ are $\frac12$.
Each $I_m$ is a Hankel transform of order $m$, and Parseval's theorem applies to each as in Part 1, Section 1,
Step 6: with $c_m(\vartheta)$ the factor in front of $J_m$ in Step 3,
$\int_0^\infty|I_m|^2r\,dr = k^{-2}\int_0^\alpha|c_m|^2(\sin\vartheta\cos\vartheta)^{-1}d\vartheta$. The three terms
add to
\begin{equation}\int_0^\infty\langle|\mathbf E|^2\rangle_\psi\,r\,dr = \frac{1}{k^2}\int_0^\alpha|A|^2\sin\vartheta\,
\big[(1 + \cos\vartheta)^2 + 2\sin^2\!\vartheta + (1 - \cos\vartheta)^2\big]\,d\vartheta
= \frac{4}{k^2}\int_0^\alpha|A|^2\sin\vartheta\,d\vartheta ,\end{equation}
since the bracket equals 4 at all $\vartheta$. This is $|C|^2 = 1/\pi^2$ times the integral of the scalar field with
the same $A$ (Part 1, Section 1, Step 6), in all planes $z$ and at all NA. The lens turns the polarization of each
ray and keeps its amplitude ($|\hat{\mathbf e}| = 1$, Section 2), so $E_x$, $E_y$ and $E_z$ share among them the
integral of the scalar field.

\chk{The engine (model \texttt{'rw'}, 750 nm, air) at $45^\circ$ to the polarization, where $|\mathbf E|^2$ equals
its mean over $\psi$, integrated as in Part 1: the integral is $4(1 - \cos\alpha)/k^2$ to $3\times10^{-4}$ at NA 0.5
and 0.9, in the focal plane and 2~$\mu$m from it.}

\section{Gauss's law}
\step{1} In a region without charges $\nabla\cdot\mathbf E = 0$. Near the focus at low NA, $E_x \approx -iI_0$, and
$E_y$ is of second order in the angle. Each plane wave in the sum varies along $z$ as
$e^{ikz\cos\vartheta} \approx e^{ikz}$, so $\partial E_z/\partial z \approx ikE_z$. Gauss's law then reads
\begin{equation}ik\,E_z \approx -\frac{\partial E_x}{\partial x} = i\,\frac{\partial I_0}{\partial x} ,\qquad
E_z \approx \frac1k\,\frac{\partial I_0}{\partial x} .\end{equation}

\step{2} With $J_0' = -J_1$ and $\partial r/\partial x = \cos\psi$,
$\partial J_0(kr\sin\vartheta)/\partial x = -k\sin\vartheta\,J_1(kr\sin\vartheta)\cos\psi$, so
\begin{equation}\frac{\partial I_0}{\partial x} = -k\cos\psi\int_0^\alpha A\sqrt{\cos\vartheta}\,\sin^2\!\vartheta\,
(1 + \cos\vartheta)\,J_1\,e^{ikz\cos\vartheta}\,d\vartheta \approx -2k\cos\psi\,I_1 ,\end{equation}
where $1 + \cos\vartheta \approx 2$ at low NA.

\step{3} Putting Step 2 into Step 1, $E_z \approx -2I_1\cos\psi$: the real factor $-2$ of Eq.~(\ref{eq:rw}).
With $-2i$ the field would not satisfy Gauss's law.

\chk{NA 0.9, air, 750 nm, on a grid of $5\times5\times5$ points within 0.6~$\mu$m of the focus, by central differences
with a step of 1 nm: the largest $|\nabla\cdot\mathbf E|$ is $1.1\times10^{-6}$ of $k|\mathbf E(0)|$ with the factor
$-2$ and 0.37 of it with $-2i$.}

\section{The low-NA limit}
\step{1} For a narrow cone, $1 + \cos\vartheta \to 2$, while $\sin^2\!\vartheta$ and $\sin\vartheta(1 - \cos\vartheta)$
are smaller than $\sin\vartheta$ by factors of order $\vartheta$ and $\vartheta^2/2$. So $I_1$ and $I_2$ vanish
against $I_0$, and $I_0$ becomes twice the scalar Debye integral of Part 1 (Section 1, Step 4) without its factor $2\pi$.

\step{2} The field becomes $\mathbf E \approx -iI_0\,\hat{\mathbf x}$: the scalar field, polarized along $x$. This
is the limit of the two-color models of [CS00], [LS02] and [LS03].

\chk{Largest share of the peak intensity $|I_0(0)|^2$ in the focal plane, 750 nm, air:
\begin{center}\small\begin{tabular}{lcccc}\toprule
& NA 0.1 & NA 0.5 & NA 0.7 & NA 0.9\\\midrule
$|E_z|^2 = 4|I_1|^2$ (along the polarization) & 0.13\% & 3.52\% & 7.62\% & 15.19\%\\
$|I_2|^2$ & 0.000\% & 0.023\% & 0.111\% & 0.467\%\\\bottomrule
\end{tabular}\end{center}
At NA 0.1 the paraxial, scalar Debye and vector profiles agree to within $3\times10^{-3}$ across and along the axis
(\texttt{tests/test\_unified\_framework.py}).}

\section{The spot across the axis depends on the direction}
\step{1, along the polarization} At $\psi = 0$: $\cos2\psi = 1$, $\sin2\psi = 0$, $\cos\psi = 1$, and
\begin{equation}|\mathbf E|^2 = |I_0 + I_2|^2 + 4|I_1|^2 .\end{equation}

\step{2, across the polarization} At $\psi = \pi/2$: $\cos2\psi = -1$, $\sin2\psi = 0$, $\cos\psi = 0$, and
\begin{equation}|\mathbf E|^2 = |I_0 - I_2|^2 ,\end{equation}
with no longitudinal field.

\step{3} In the focal plane near the axis $I_0$ is positive and falls with $r$, and $I_2$ is positive and grows as
$r^2$, because $J_2(v) \approx v^2/8$. Adding $I_2$ makes the intensity fall more slowly and subtracting it more
quickly; $4|I_1|^2$ is zero on the axis and largest beside it. The spot is therefore wider along the polarization and
narrower across it.

\chk{FWHM in the focal plane, one beam, 750 nm, air ($\mu$m):
\begin{center}\small\begin{tabular}{lcccc}\toprule
& NA 0.1 & NA 0.5 & NA 0.7 & NA 0.9\\\midrule
along the polarization & 3.8659 & 0.8131 & 0.6199 & 0.5481\\
paraxial, $0.5145\,\lambda/\mathrm{NA}$ & 3.8587 & 0.7717 & 0.5512 & 0.4287\\
across the polarization & 3.8570 & 0.7623 & 0.5363 & 0.4042\\\bottomrule
\end{tabular}\end{center}
The paraxial spot lies between the two at each NA. Across the polarization the vector spot is 1.2\% narrower than
the paraxial one at NA 0.5 and 5.7\% narrower at NA 0.9; along it, 5.4\% and 28\% wider. The depth of focus changes
more: the paraxial model makes it 7\% too long at NA 0.5 and 36\% too long at NA 0.9 (Part 1, Section 7).}

\section{The plane that the current work computes}
\step{1} Both beams are polarized along the lab $y$ axis, perpendicular to the plane of the two lens axes (the
$xz$-plane). In the frame of either beam, a point of that plane lies at $\pm\pi/2$ from the polarization, so all
points of the plane are across the polarization.

\step{2} On that plane each beam therefore has $|\mathbf E|^2 = |I_0 - I_2|^2$, its field points along $y$, and
$E_z = 0$. The stored maps and widths in the plane of the lens axes are across the polarization. This is why the 2CE
spot of the vector model is 1.2\% narrower than the paraxial one at NA 0.5 (0.510 against 0.516~$\mu$m) and 5.6\%
narrower at NA 0.9 (0.270 against 0.286~$\mu$m).

\step{3} On the axis $J_1(0) = J_2(0) = 0$, so $I_1 = I_2 = 0$ and $|\mathbf E|^2 = |I_0|^2$. The depth of focus is
the Fourier transform of the weight $\sqrt t\,(1 + t)$ of Part 1, Section 7, and does not depend on the polarization.

\chk{The two-beam engine gives the FWHM of [Rom03], 472.8 nm (NA 0.8, $n = 1.3$, 750 nm, across the polarization), and
its vector field equals that of the Romallosa replication (\texttt{mcdo/rw\_integrals.py}) to $10^{-9}$ for both
polarizations and along the axis (\texttt{tests/test\_unified\_framework.py}).}

\step{4, what the factor of $E_z$ changes} $|E_z|^2$ is the same with $-2$ and with $-2i$, so the intensity of one
beam, the 2CE signal $I_1I_2$ and the Romallosa replication do not depend on it. The factor matters where the $E_z$ of
one beam adds to a component of another beam before the square is taken. That happens in 2PE with crossed lenses when
the polarization lies in the plane of the two beams: the $E_z$ of beam 2 then points along the lab $x$ axis and adds to
the $E_x$ of beam 1. In the stored high-NA pupil study 208 rows of that case changed when the factor was corrected;
none of them is on a slide. With the polarization perpendicular to the plane, as in all the stored maps and in the
pulse and delay runs, $E_z = 0$ on the plane (Step 2).

\section*{References}
\small
[CS00] M. O. Cambaliza and C. Saloma, Opt.\ Commun.\ 184, 25--35 (2000).\\
\relax[LS02] M. Lim and C. Saloma, Opt.\ Commun.\ 207, 121--130 (2002).\\
\relax[LS03] M. Lim and C. Saloma, Appl.\ Opt.\ 42(17), 3398--3406 (2003).\\
\relax[Rom03] K. M. Romallosa, J. Bantang and C. Saloma, Phys.\ Rev.\ A 68, 033812 (2003).\\
B. Richards and E. Wolf, Proc.\ R.\ Soc.\ London Ser.\ A 253, 358--379 (1959).\\
L. Novotny and B. Hecht, \emph{Principles of Nano-Optics}, Cambridge University Press (2006), Ch.~3.
\end{document}
